\begin{table}[ht]
\centering
\caption{Prior distributions for the DDM and OUM parameters.}
\label{tab:priors}
\begin{tabular}{ll cc}
\toprule
 & Parameter & DDM & OUM \\
\midrule
\multicolumn{4}{l}{\textit{Study 1 --- fast tasks}} \\
 & Drift rate $\nu$                 & $\Gamma(3.5,\,1.0)$               & $\Gamma(3.5,\,1.0)$ \\
 & Boundary separation $a$          & $8\,\sigma(\mathcal{N}(-1.6,\,1))$ & $8\,\sigma(\mathcal{N}(-0.25,\,1))$ \\
 & Non-decision time $\tau$         & $\mathcal{U}(0.1,\,1.0)$          & $\mathcal{U}(0.1,\,1.0)$ \\
 & Self-excitation $k$              & ---                               & $\Gamma(8.0,\,0.5)$ \\
\addlinespace
\multicolumn{4}{l}{\textit{Study 1 --- slow tasks}} \\
 & Drift rate $\nu$                 & $\Gamma(6.0,\,0.25)$              & $\Gamma(6.0,\,0.25)$ \\
 & Boundary separation $a$          & $10\,\sigma(\mathcal{N}(-0.7,\,1))$ & $10\,\sigma(\mathcal{N}(0.5,\,1.2))$ \\
 & Non-decision time $\tau$         & $\mathcal{U}(0.1,\,3.0)$          & $\mathcal{U}(0.1,\,3.0)$ \\
 & Self-excitation $k$              & ---                               & $\Gamma(3.0,\,0.3)$ \\
\addlinespace
\multicolumn{4}{l}{\textit{Study 2 --- Race IAT}} \\
 & Drift rates $\nu_1,\,\nu_2$      & $\Gamma(3.5,\,1.0)$               & $\Gamma(3.5,\,1.0)$ \\
 & Boundary separations $a_1,\,a_2$ & $8\,\sigma(\mathcal{N}(-1.6,\,1))$ & $8\,\sigma(\mathcal{N}(-0.25,\,1))$ \\
 & Non-decision time, correct $\tau_c$ & $\mathcal{U}(0.1,\,1.0)$       & $\mathcal{U}(0.1,\,1.0)$ \\
 & Non-decision time, error $\tau_e$   & $\Gamma(2.0,\,0.3)$           & $\Gamma(2.0,\,0.3)$ \\
 & Self-excitation $k$              & ---                               & $\Gamma(8.0,\,0.5)$ \\
\bottomrule
\end{tabular}

\vspace{4pt}
\noindent\textit{Note.} $\sigma(\cdot)$ denotes the logistic sigmoid, giving a smoothly bounded prior on $(0, L)$ for the boundary separation (with $L = 8$ for fast tasks and the IAT, $L = 10$ for slow tasks). Gamma priors use the shape--scale parameterization, $\Gamma(\alpha, \theta)$ with mean $\alpha\theta$. The OUM uses a higher-centered boundary prior than the DDM because OUM boundary estimates tend to be larger. The Study 2 (IAT) priors match the Study 1 fast-task priors, with an additional error non-decision time $\tau_e$ that absorbs the extra response-correction time recorded for IAT errors. $k$ is restricted to positive values (self-excitation); for the slow tasks its prior places substantial mass near zero, where the OUM reduces to the DDM.
\end{table}
